\chapter{Evaluation}
\label{chap:eval}

\begin{itemize}
	\item Introduce the benchmarking framework.
	\item Goal: show where the framework is strong and where it has limitations.
	\item Interpret the results from the point of view of the states: what the states give us and what not.
	\item Short notes in the sections and subsections, referenced later in Chapters~\ref{chap:discussion} and~\ref{chap:conclusion}.
\end{itemize}

\section{Padella Benchmark}
\label{sec:eval:padella}
\begin{itemize}
	\item Introduce the Padella benchmark.
\end{itemize}

\begin{figure}[tb]
	\centering
	\begin{tikzpicture}[font=\small, >={Stealth},
			box/.style={draw, rounded corners=2pt, fill=black!6, minimum width=2.4cm, minimum height=0.7cm}]
		\node[box] (cf) at (0, 0) {\texttt{cf}};
		\node[box] (fast) at (0, -1.0) {\texttt{time\_fast}};
		\node[box] (slow) at (0, -2.0) {\texttt{time\_slow}};
		\node[box] (intra) at (5, 0) {Intra-case};
		\node[box] (resource) at (5, -1.0) {Resource};
		\node[box] (inter) at (5, -2.0) {Inter-case};
		\draw[->] (cf) -- (intra);
		\draw (fast.east) -- (2.2, -1.0) -- (2.2, -2.0);
		\draw (slow.east) -- (2.2, -2.0);
		\draw[->] (2.2, -1.0) -- (resource.west);
		\draw[->, dashed] (2.2, -2.0) -- (inter.west);
	\end{tikzpicture}
	\caption[Padella drift types and our perspectives]{%
		Perspectives that capture the drift types of the Padella benchmark. Dashed: also responds.}
	\label{fig:eval:padella:perspectives}
\end{figure}

\begin{itemize}
	\item Global results over all logs of the benchmark first.
	\item For three thresholds: confusion matrix (TP, FP, TN, FN), precision, recall and F1.
	\item Then the ROC curve with its AUC.
	\item Note: strengths, limitations and what the states show here, for Chapters~\ref{chap:discussion} and~\ref{chap:conclusion}.
\end{itemize}

% mock data, to be replaced by the results
\begin{figure}[tb]
	\centering
	\def\confmat#1#2#3#4{%
		\begin{tikzpicture}[font=\footnotesize, cell/.style={draw, minimum width=1.2cm, minimum height=0.7cm}]
			\node[cell, fill=black!15] at (0, 0) {#1};
			\node[cell] at (1.2, 0) {#2};
			\node[cell] at (0, -0.7) {#3};
			\node[cell, fill=black!15] at (1.2, -0.7) {#4};
			\node[font=\scriptsize] at (0, 0.55) {drift};
			\node[font=\scriptsize] at (1.2, 0.55) {no drift};
			\node[font=\scriptsize, anchor=east] at (-0.65, 0) {drift};
			\node[font=\scriptsize, anchor=east] at (-0.65, -0.7) {no drift};
		\end{tikzpicture}}
	\begin{subfigure}[b]{0.32\textwidth}
		\centering
		\confmat{41}{4}{58}{897}
		\caption{$\tau_1$}
	\end{subfigure}
	\hfill
	\begin{subfigure}[b]{0.32\textwidth}
		\centering
		\confmat{36}{9}{21}{934}
		\caption{$\tau_2$}
	\end{subfigure}
	\hfill
	\begin{subfigure}[b]{0.32\textwidth}
		\centering
		\confmat{27}{18}{6}{949}
		\caption{$\tau_3$}
	\end{subfigure}
	\caption[Confusion matrices on the Padella benchmark (mock data)]{%
		Confusion matrices over all Padella logs for three thresholds $\tau_1 < \tau_2 < \tau_3$. Rows: actual, columns: detected (mock data).}
	\label{fig:eval:padella:confusion}
\end{figure}

\begin{table}[tb]
	\centering
	\caption[Precision, recall and F1 on the Padella benchmark (mock data)]{%
		Precision, recall and F1 over all Padella logs for the thresholds of Figure~\ref{fig:eval:padella:confusion} (mock data).}
	\label{tab:eval:padella:metrics}
	\small
	\begin{tabular}{@{}l|r|r|r@{}}
		\hline
		Threshold & Precision & Recall & F1    \\
		\hline
		$\tau_1$  & 0.414     & 0.911  & 0.569 \\
		\hline
		$\tau_2$  & 0.632     & 0.800  & 0.706 \\
		\hline
		$\tau_3$  & 0.818     & 0.600  & 0.692 \\
		\hline
	\end{tabular}
\end{table}

\begin{figure}[tb]
	\centering
	\begin{tikzpicture}[x=4cm, y=4cm, font=\footnotesize, >={Stealth}]
		\draw[->] (0, 0) -- (1.08, 0);
		\draw[->] (0, 0) -- (0, 1.08);
		\node at (0.5, -0.15) {false positive rate};
		\node[rotate=90] at (-0.24, 0.5) {true positive rate};
		\foreach \t in {0, 0.5, 1} {
			\draw (\t, 0) -- (\t, -0.02) node[below] {\t};
			\draw (0, \t) -- (-0.02, \t) node[left] {\t};
		}
		\draw[dashed, black!40] (0, 0) -- (1, 1);
		\draw[thick] (0, 0) -- (0.006, 0.6) -- (0.022, 0.8) -- (0.061, 0.911) -- (0.2, 0.96) -- (0.5, 0.99) -- (1, 1);
		\foreach \x/\y/\l in {0.006/0.6/3, 0.022/0.8/2, 0.061/0.911/1} {
			\fill (\x, \y) circle (0.012);
			\node[anchor=west] at (\x + 0.02, \y) {$\tau_\l$};
		}
		\node[anchor=south east] at (1, 0.05) {AUC $= 0.95$};
	\end{tikzpicture}
	\caption[ROC curve on the Padella benchmark (mock data)]{%
		ROC curve over all Padella logs with the three thresholds of Figure~\ref{fig:eval:padella:confusion} (mock data).}
	\label{fig:eval:padella:roc}
\end{figure}

\subsection{Control-Flow Drifts}
\label{sec:eval:padella:cf}
\begin{itemize}
	\item Results of the logs with control-flow drifts (\texttt{cf}).
	\item How well the framework detects this drift type.
	\item For three thresholds: confusion matrix, precision, recall and F1, then the ROC curve with its AUC.
	\item Note: strengths, limitations and what the states show here, for Chapters~\ref{chap:discussion} and~\ref{chap:conclusion}.
\end{itemize}

% mock data, to be replaced by the results
\begin{figure}[tb]
	\centering
	\begin{tikzpicture}[x=4cm, y=4cm, font=\footnotesize, >={Stealth}]
		\draw[->] (0, 0) -- (1.08, 0);
		\draw[->] (0, 0) -- (0, 1.08);
		\node at (0.5, -0.15) {false positive rate};
		\node[rotate=90] at (-0.24, 0.5) {true positive rate};
		\foreach \t in {0, 0.5, 1} {
			\draw (\t, 0) -- (\t, -0.02) node[below] {\t};
			\draw (0, \t) -- (-0.02, \t) node[left] {\t};
		}
		\draw[dashed, black!40] (0, 0) -- (1, 1);
		\draw[thick] (0, 0) -- (0.01, 0.73) -- (0.04, 0.93) -- (0.15, 1) -- (1, 1);
		\draw[thick, densely dashed] (0, 0) -- (0.03, 0.47) -- (0.09, 0.73) -- (0.3, 0.93) -- (1, 1);
		\foreach \y/\style/\l in {0.3/solid/{sudden, AUC $= 0.98$}, 0.21/densely dashed/{gradual, AUC $= 0.92$}} {
			\draw[thick, \style] (0.42, \y) -- (0.5, \y);
			\node[anchor=west] at (0.51, \y) {\l};
		}
	\end{tikzpicture}
	\caption[ROC curves for control-flow drifts (mock data)]{%
		ROC curves of the Padella logs with control-flow drifts, per drift mode (mock data).}
	\label{fig:eval:padella:cf_roc}
\end{figure}

\subsection{Fast Time Drifts}
\label{sec:eval:padella:time_fast}
\begin{itemize}
	\item Results of the logs with fast time drifts (\texttt{time\_fast}).
	\item How well the framework detects this drift type.
	\item For three thresholds: confusion matrix, precision, recall and F1, then the ROC curve with its AUC.
	\item Note: strengths, limitations and what the states show here, for Chapters~\ref{chap:discussion} and~\ref{chap:conclusion}.
\end{itemize}

\subsection{Slow Time Drifts}
\label{sec:eval:padella:time_slow}
\begin{itemize}
	\item Results of the logs with slow time drifts (\texttt{time\_slow}).
	\item How well the framework detects this drift type.
	\item For three thresholds: confusion matrix, precision, recall and F1, then the ROC curve with its AUC.
	\item Note: strengths, limitations and what the states show here, for Chapters~\ref{chap:discussion} and~\ref{chap:conclusion}.
\end{itemize}

\subsection{Sudden Drifts}
\label{sec:eval:padella:sudden}
\begin{itemize}
	\item Results of the logs with sudden drifts, over all drift types.
	\item How well the framework detects this drift mode.
	\item For three thresholds: confusion matrix, precision, recall and F1, then the ROC curve with its AUC.
	\item Note: strengths, limitations and what the states show here, for Chapters~\ref{chap:discussion} and~\ref{chap:conclusion}.
\end{itemize}

\subsection{Gradual Drifts}
\label{sec:eval:padella:gradual}
\begin{itemize}
	\item Results of the logs with gradual drifts, over all drift types.
	\item How well the framework detects this drift mode.
	\item For three thresholds: confusion matrix, precision, recall and F1, then the ROC curve with its AUC.
	\item Note: strengths, limitations and what the states show here, for Chapters~\ref{chap:discussion} and~\ref{chap:conclusion}.
\end{itemize}

\section{Rheon Benchmark}
\label{sec:eval:rheon}
\begin{itemize}
	\item Three sets of Rheon logs: single-drift, multi-drift and mixed-drift logs.
	\item In each set, the global results first, then the results per drift perspective and per drift mode.
	\item Note: strengths, limitations and what the states show here, for Chapters~\ref{chap:discussion} and~\ref{chap:conclusion}.
\end{itemize}

\subsection{Single-Drift Logs}
\label{sec:eval:rheon:single}
\begin{itemize}
	\item Logs with one drift each.
	\item Global results over all single-drift logs first.
	\item For three thresholds: confusion matrix, precision, recall and F1, then the ROC curve with its AUC.
	\item Note: strengths, limitations and what the states show here, for Chapters~\ref{chap:discussion} and~\ref{chap:conclusion}.
\end{itemize}

% mock data, to be replaced by the results
\begin{table}[tb]
	\centering
	\caption[Results on the single-drift Rheon logs (mock data)]{%
		Precision, recall, F1 and AUC of the single-drift Rheon logs per drift perspective for the middle threshold $\tau_2$ (mock data).}
	\label{tab:eval:rheon:single}
	\small
	\begin{tabular}{@{}l|r|r|r|r@{}}
		\hline
		Drift perspective & Precision & Recall & F1    & AUC  \\
		\hline
		Intra-case        & 0.880     & 0.917  & 0.898 & 0.97 \\
		\hline
		Resource          & 0.714     & 0.833  & 0.769 & 0.91 \\
		\hline
		Inter-case        & 0.650     & 0.722  & 0.684 & 0.86 \\
		\hline
	\end{tabular}
\end{table}

\subsubsection{Intra-Case Drifts}
\label{sec:eval:rheon:single:intra}
\begin{itemize}
	\item Results of the single-drift logs with an intra-case drift (\texttt{control\_flow}).
	\item For three thresholds: confusion matrix, precision, recall and F1, then the ROC curve with its AUC.
	\item Note: strengths, limitations and what the states show here, for Chapters~\ref{chap:discussion} and~\ref{chap:conclusion}.
\end{itemize}

\subsubsection{Resource Drifts}
\label{sec:eval:rheon:single:resource}
\begin{itemize}
	\item Results of the single-drift logs with a resource drift.
	\item For three thresholds: confusion matrix, precision, recall and F1, then the ROC curve with its AUC.
	\item Note: strengths, limitations and what the states show here, for Chapters~\ref{chap:discussion} and~\ref{chap:conclusion}.
\end{itemize}

\subsubsection{Inter-Case Drifts}
\label{sec:eval:rheon:single:inter}
\begin{itemize}
	\item Results of the single-drift logs with an inter-case drift.
	\item For three thresholds: confusion matrix, precision, recall and F1, then the ROC curve with its AUC.
	\item Note: strengths, limitations and what the states show here, for Chapters~\ref{chap:discussion} and~\ref{chap:conclusion}.
\end{itemize}

\subsubsection{Sudden Drifts}
\label{sec:eval:rheon:single:sudden}
\begin{itemize}
	\item Results of the single-drift logs with a sudden drift.
	\item For three thresholds: confusion matrix, precision, recall and F1, then the ROC curve with its AUC.
	\item Note: strengths, limitations and what the states show here, for Chapters~\ref{chap:discussion} and~\ref{chap:conclusion}.
\end{itemize}

\subsubsection{Gradual Drifts}
\label{sec:eval:rheon:single:gradual}
\begin{itemize}
	\item Results of the single-drift logs with a gradual drift.
	\item For three thresholds: confusion matrix, precision, recall and F1, then the ROC curve with its AUC.
	\item Note: strengths, limitations and what the states show here, for Chapters~\ref{chap:discussion} and~\ref{chap:conclusion}.
\end{itemize}

\subsection{Multi-Drift Logs}
\label{sec:eval:rheon:multi}
\begin{itemize}
	\item Logs with several drifts, all from the same perspective.
	\item Global results over all multi-drift logs first.
	\item For three thresholds: confusion matrix, precision, recall and F1, then the ROC curve with its AUC.
	\item Note: strengths, limitations and what the states show here, for Chapters~\ref{chap:discussion} and~\ref{chap:conclusion}.
\end{itemize}

\subsubsection{Intra-Case Drifts}
\label{sec:eval:rheon:multi:intra}
\begin{itemize}
	\item Results of the intra-case drifts in the multi-drift logs.
	\item For three thresholds: confusion matrix, precision, recall and F1, then the ROC curve with its AUC.
	\item Note: strengths, limitations and what the states show here, for Chapters~\ref{chap:discussion} and~\ref{chap:conclusion}.
\end{itemize}

\subsubsection{Resource Drifts}
\label{sec:eval:rheon:multi:resource}
\begin{itemize}
	\item Results of the resource drifts in the multi-drift logs.
	\item For three thresholds: confusion matrix, precision, recall and F1, then the ROC curve with its AUC.
	\item Note: strengths, limitations and what the states show here, for Chapters~\ref{chap:discussion} and~\ref{chap:conclusion}.
\end{itemize}

\subsubsection{Inter-Case Drifts}
\label{sec:eval:rheon:multi:inter}
\begin{itemize}
	\item Results of the inter-case drifts in the multi-drift logs.
	\item For three thresholds: confusion matrix, precision, recall and F1, then the ROC curve with its AUC.
	\item Note: strengths, limitations and what the states show here, for Chapters~\ref{chap:discussion} and~\ref{chap:conclusion}.
\end{itemize}

\subsubsection{Sudden Drifts}
\label{sec:eval:rheon:multi:sudden}
\begin{itemize}
	\item Results of the sudden drifts in the multi-drift logs.
	\item For three thresholds: confusion matrix, precision, recall and F1, then the ROC curve with its AUC.
	\item Note: strengths, limitations and what the states show here, for Chapters~\ref{chap:discussion} and~\ref{chap:conclusion}.
\end{itemize}

\subsubsection{Gradual Drifts}
\label{sec:eval:rheon:multi:gradual}
\begin{itemize}
	\item Results of the gradual drifts in the multi-drift logs.
	\item For three thresholds: confusion matrix, precision, recall and F1, then the ROC curve with its AUC.
	\item Note: strengths, limitations and what the states show here, for Chapters~\ref{chap:discussion} and~\ref{chap:conclusion}.
\end{itemize}

\subsection{Mixed-Drift Logs}
\label{sec:eval:rheon:mixed}
\begin{itemize}
	\item Logs with several drifts from different perspectives.
	\item Global results over all mixed-drift logs first.
	\item For three thresholds: confusion matrix, precision, recall and F1, then the ROC curve with its AUC.
	\item Note: strengths, limitations and what the states show here, for Chapters~\ref{chap:discussion} and~\ref{chap:conclusion}.
\end{itemize}

\subsubsection{Intra-Case Drifts}
\label{sec:eval:rheon:mixed:intra}
\begin{itemize}
	\item Results of the intra-case drifts in the mixed-drift logs.
	\item For three thresholds: confusion matrix, precision, recall and F1, then the ROC curve with its AUC.
	\item Note: strengths, limitations and what the states show here, for Chapters~\ref{chap:discussion} and~\ref{chap:conclusion}.
\end{itemize}

\subsubsection{Resource Drifts}
\label{sec:eval:rheon:mixed:resource}
\begin{itemize}
	\item Results of the resource drifts in the mixed-drift logs.
	\item For three thresholds: confusion matrix, precision, recall and F1, then the ROC curve with its AUC.
	\item Note: strengths, limitations and what the states show here, for Chapters~\ref{chap:discussion} and~\ref{chap:conclusion}.
\end{itemize}

\subsubsection{Inter-Case Drifts}
\label{sec:eval:rheon:mixed:inter}
\begin{itemize}
	\item Results of the inter-case drifts in the mixed-drift logs.
	\item For three thresholds: confusion matrix, precision, recall and F1, then the ROC curve with its AUC.
	\item Note: strengths, limitations and what the states show here, for Chapters~\ref{chap:discussion} and~\ref{chap:conclusion}.
\end{itemize}

\subsubsection{Sudden Drifts}
\label{sec:eval:rheon:mixed:sudden}
\begin{itemize}
	\item Results of the sudden drifts in the mixed-drift logs.
	\item For three thresholds: confusion matrix, precision, recall and F1, then the ROC curve with its AUC.
	\item Note: strengths, limitations and what the states show here, for Chapters~\ref{chap:discussion} and~\ref{chap:conclusion}.
\end{itemize}

\subsubsection{Gradual Drifts}
\label{sec:eval:rheon:mixed:gradual}
\begin{itemize}
	\item Results of the gradual drifts in the mixed-drift logs.
	\item For three thresholds: confusion matrix, precision, recall and F1, then the ROC curve with its AUC.
	\item Note: strengths, limitations and what the states show here, for Chapters~\ref{chap:discussion} and~\ref{chap:conclusion}.
\end{itemize}
